% Options for packages loaded elsewhere
% Options for packages loaded elsewhere
\PassOptionsToPackage{unicode}{hyperref}
\PassOptionsToPackage{hyphens}{url}
\PassOptionsToPackage{dvipsnames,svgnames,x11names}{xcolor}
%
\documentclass[
  letterpaper,
  DIV=11,
  numbers=noendperiod]{scrartcl}
\usepackage{xcolor}
\usepackage{amsmath,amssymb}
\setcounter{secnumdepth}{5}
\usepackage{iftex}
\ifPDFTeX
  \usepackage[T1]{fontenc}
  \usepackage[utf8]{inputenc}
  \usepackage{textcomp} % provide euro and other symbols
\else % if luatex or xetex
  \usepackage{unicode-math} % this also loads fontspec
  \defaultfontfeatures{Scale=MatchLowercase}
  \defaultfontfeatures[\rmfamily]{Ligatures=TeX,Scale=1}
\fi
\usepackage{lmodern}
\ifPDFTeX\else
  % xetex/luatex font selection
\fi
% Use upquote if available, for straight quotes in verbatim environments
\IfFileExists{upquote.sty}{\usepackage{upquote}}{}
\IfFileExists{microtype.sty}{% use microtype if available
  \usepackage[]{microtype}
  \UseMicrotypeSet[protrusion]{basicmath} % disable protrusion for tt fonts
}{}
\makeatletter
\@ifundefined{KOMAClassName}{% if non-KOMA class
  \IfFileExists{parskip.sty}{%
    \usepackage{parskip}
  }{% else
    \setlength{\parindent}{0pt}
    \setlength{\parskip}{6pt plus 2pt minus 1pt}}
}{% if KOMA class
  \KOMAoptions{parskip=half}}
\makeatother
% Make \paragraph and \subparagraph free-standing
\makeatletter
\ifx\paragraph\undefined\else
  \let\oldparagraph\paragraph
  \renewcommand{\paragraph}{
    \@ifstar
      \xxxParagraphStar
      \xxxParagraphNoStar
  }
  \newcommand{\xxxParagraphStar}[1]{\oldparagraph*{#1}\mbox{}}
  \newcommand{\xxxParagraphNoStar}[1]{\oldparagraph{#1}\mbox{}}
\fi
\ifx\subparagraph\undefined\else
  \let\oldsubparagraph\subparagraph
  \renewcommand{\subparagraph}{
    \@ifstar
      \xxxSubParagraphStar
      \xxxSubParagraphNoStar
  }
  \newcommand{\xxxSubParagraphStar}[1]{\oldsubparagraph*{#1}\mbox{}}
  \newcommand{\xxxSubParagraphNoStar}[1]{\oldsubparagraph{#1}\mbox{}}
\fi
\makeatother


\usepackage{longtable,booktabs,array}
\usepackage{calc} % for calculating minipage widths
% Correct order of tables after \paragraph or \subparagraph
\usepackage{etoolbox}
\makeatletter
\patchcmd\longtable{\par}{\if@noskipsec\mbox{}\fi\par}{}{}
\makeatother
% Allow footnotes in longtable head/foot
\IfFileExists{footnotehyper.sty}{\usepackage{footnotehyper}}{\usepackage{footnote}}
\makesavenoteenv{longtable}
\usepackage{graphicx}
\makeatletter
\newsavebox\pandoc@box
\newcommand*\pandocbounded[1]{% scales image to fit in text height/width
  \sbox\pandoc@box{#1}%
  \Gscale@div\@tempa{\textheight}{\dimexpr\ht\pandoc@box+\dp\pandoc@box\relax}%
  \Gscale@div\@tempb{\linewidth}{\wd\pandoc@box}%
  \ifdim\@tempb\p@<\@tempa\p@\let\@tempa\@tempb\fi% select the smaller of both
  \ifdim\@tempa\p@<\p@\scalebox{\@tempa}{\usebox\pandoc@box}%
  \else\usebox{\pandoc@box}%
  \fi%
}
% Set default figure placement to htbp
\def\fps@figure{htbp}
\makeatother





\setlength{\emergencystretch}{3em} % prevent overfull lines

\providecommand{\tightlist}{%
  \setlength{\itemsep}{0pt}\setlength{\parskip}{0pt}}



 


\KOMAoption{captions}{tableheading}
\makeatletter
\@ifpackageloaded{tcolorbox}{}{\usepackage[skins,breakable]{tcolorbox}}
\@ifpackageloaded{fontawesome5}{}{\usepackage{fontawesome5}}
\definecolor{quarto-callout-color}{HTML}{909090}
\definecolor{quarto-callout-note-color}{HTML}{0758E5}
\definecolor{quarto-callout-important-color}{HTML}{CC1914}
\definecolor{quarto-callout-warning-color}{HTML}{EB9113}
\definecolor{quarto-callout-tip-color}{HTML}{00A047}
\definecolor{quarto-callout-caution-color}{HTML}{FC5300}
\definecolor{quarto-callout-color-frame}{HTML}{acacac}
\definecolor{quarto-callout-note-color-frame}{HTML}{4582ec}
\definecolor{quarto-callout-important-color-frame}{HTML}{d9534f}
\definecolor{quarto-callout-warning-color-frame}{HTML}{f0ad4e}
\definecolor{quarto-callout-tip-color-frame}{HTML}{02b875}
\definecolor{quarto-callout-caution-color-frame}{HTML}{fd7e14}
\makeatother
\makeatletter
\@ifpackageloaded{caption}{}{\usepackage{caption}}
\AtBeginDocument{%
\ifdefined\contentsname
  \renewcommand*\contentsname{Table of contents}
\else
  \newcommand\contentsname{Table of contents}
\fi
\ifdefined\listfigurename
  \renewcommand*\listfigurename{List of Figures}
\else
  \newcommand\listfigurename{List of Figures}
\fi
\ifdefined\listtablename
  \renewcommand*\listtablename{List of Tables}
\else
  \newcommand\listtablename{List of Tables}
\fi
\ifdefined\figurename
  \renewcommand*\figurename{Figure}
\else
  \newcommand\figurename{Figure}
\fi
\ifdefined\tablename
  \renewcommand*\tablename{Table}
\else
  \newcommand\tablename{Table}
\fi
}
\@ifpackageloaded{float}{}{\usepackage{float}}
\floatstyle{ruled}
\@ifundefined{c@chapter}{\newfloat{codelisting}{h}{lop}}{\newfloat{codelisting}{h}{lop}[chapter]}
\floatname{codelisting}{Listing}
\newcommand*\listoflistings{\listof{codelisting}{List of Listings}}
\makeatother
\makeatletter
\makeatother
\makeatletter
\@ifpackageloaded{caption}{}{\usepackage{caption}}
\@ifpackageloaded{subcaption}{}{\usepackage{subcaption}}
\makeatother
\makeatletter
\@ifpackageloaded{tcolorbox}{}{\usepackage[many]{tcolorbox}}
\makeatother
\renewcommand{\ul}[1]{\underline{#1}} % because soul conflicts with tcolorbox
%%%% ---simpletextbox preamble ----- %%%%%
% arguments: #1 background color
\newtcolorbox{simpletextbox}[1]{breakable, sharp corners, boxrule=0pt, 
    boxsep=1em, beforeafter skip=1em, leftright skip=1em, 
    colback=#1-color1}%


%\renewcommand{\ul}[1]{unterstrich #1}

%%%% --- end simpletextbox preamble ----- %%%%%
%%%% ---foldboxy preamble ----- %%%%%

\definecolor{fobx-default-color1}{HTML}{c7c7d0}
\definecolor{fobx-default-color2}{HTML}{a3a3aa}

\definecolor{fobox-color1}{HTML}{c7c7d0}
\definecolor{fobox-color2}{HTML}{a3a3aa}

% arguments: #1 typelabelnummer: #2 titel: #3
\newenvironment{fobx}[2]{\begin{tcolorbox}[enhanced, breakable,%
attach boxed title to top*={xshift=1.4pt},
boxed title style={boxrule=0.0mm, fuzzy shadow={1pt}{-1pt}{0mm}{0.1mm}{gray}, arc=.3em, rounded corners=east, sharp corners=west}, colframe=#1-color2, colbacktitle=#1-color1, colback = white, coltitle=black,  titlerule=0mm, toprule=0pt, bottomrule=.7pt, leftrule=.3em, rightrule=0pt, outer arc=.3em,  arc=0pt,	 sharp corners = east, left=.5em, bottomtitle=1mm, toptitle=1mm,title={#2}]}
{\end{tcolorbox}}

% boxed environment with right border
\newenvironment{fobxSimple}[2]{\begin{tcolorbox}[enhanced, breakable,%
attach boxed title to top*={xshift=1.4pt},
boxed title style={boxrule=0.0mm, fuzzy shadow={1pt}{-1pt}{0mm}{0.1mm}{gray}, arc=.3em, rounded corners=east, sharp corners=west}, colframe=#1-color2, colbacktitle=#1-color1, colback = white, coltitle=black,  titlerule=0mm, toprule=0pt, bottomrule=.7pt, leftrule=.3em, rightrule=.7pt, outer arc=.3em,  	left=.5em, right=.5em, bottomtitle=1mm, toptitle=1mm,title={#2}]}
{\end{tcolorbox}}

%%%% --- end foldboxy preamble ----- %%%%%
%%==== colors from yaml ===%
\definecolor{TODO-color1}{HTML}{e7b1b4}
\definecolor{TODO-color2}{HTML}{8c3236}
\definecolor{Conjecture-color1}{HTML}{c0c0c0}
\definecolor{Conjecture-color2}{HTML}{808080}
\definecolor{Tip-color1}{HTML}{c0c0c0}
\definecolor{Tip-color2}{HTML}{808080}
\definecolor{Corollary-color1}{HTML}{c0c0c0}
\definecolor{Corollary-color2}{HTML}{808080}
\definecolor{Definition-color1}{HTML}{d999d3}
\definecolor{Definition-color2}{HTML}{a01793}
\definecolor{Feature-color1}{HTML}{c0c0c0}
\definecolor{Feature-color2}{HTML}{808080}
\definecolor{NewFeature-color1}{HTML}{cce7b1}
\definecolor{NewFeature-color2}{HTML}{86b754}
\definecolor{Theorem-color1}{HTML}{c0c0c0}
\definecolor{Theorem-color2}{HTML}{808080}
%=============%
\usepackage{bookmark}
\IfFileExists{xurl.sty}{\usepackage{xurl}}{} % add URL line breaks if available
\urlstyle{same}
\hypersetup{
  pdftitle={Custom-numbered-blocks Example},
  colorlinks=true,
  linkcolor={blue},
  filecolor={Maroon},
  citecolor={Blue},
  urlcolor={Blue},
  pdfcreator={LaTeX via pandoc}}


\title{Custom-numbered-blocks Example}
\author{}
\date{}
\begin{document}
\maketitle


\section{Custom blocks and
crossreferencing}\label{custom-blocks-and-crossreferencing}

With this filter, you can define custom div classes (environments) that
come with numbering, such as theorems, examples, exercises. The filter
currently supports output formats pdf and html.

It is possible to define shared counters for multiple classes of blocks,
similar to LaTeX's amsthm, by defining \texttt{groups}.

\phantomsection\label{cnb-1-NewFeature}
\begin{fobx}{NewFeature}{\textbf{New feature 1.1: }New options for block appearance}

The default appearance for custom divs is \textbf{foldbox}: a
collapsible similar to quarto's callouts, with a collapse button at the
bottom that makes it easier collapse long boxes, and box open to the
right. It comes with the variant \texttt{foldbox.simple}, with closed
box and no additional close button.

Two new ways a custom block can be rendered are available now:

\begin{itemize}
\tightlist
\item
  \textbf{quartothmlike} mimics the appearance of theorem blocks in
  quarto, however they look alike in both \texttt{html} and \texttt{pdf}
  formats
\item
  \textbf{textbox} a colored text box
\end{itemize}

\end{fobx}

\phantomsection\label{numbering-cross-reference-overhaul-with-slightly-changed-behaviour}
\begin{fobx}{Feature}{\textbf{Feature 1.2: }Numbering: cross reference
overhaul with slightly changed behaviour}

Numbering is by default within document and without section prefix for
single documents, or within chapter and with chapter prefix for books.
Grouped classes share the same counter, and the same default appearance.

The default behaviour can be changed by setting the item
\texttt{chapters} in yaml-key \texttt{crossreference} to \texttt{true}
or \texttt{false}, to comply with quarto's (no longer so) new
crossreferencing options.

In the present single document, yaml contains

\begin{verbatim}
crossref:
  chapters: true
\end{verbatim}

Numbered custom blocks can be cross-referenced with
\texttt{\textbackslash{}ref}.

Default numbering can be switched off for the whole class by setting the
\texttt{numbered} to false, or for an individual block by adding the
class \texttt{unnumbered}.

Crossreferences may need a re-run to update.

\phantomsection\label{nested-blocks-may-be-numbered}
\begin{fobx}{NewFeature}{\textbf{New feature: }Nested blocks may be
numbered}

With the new version, nested blocks are numbered. This does not always
make sense -- you can and probably should alter this behaviour by
specifying class \texttt{.unnumbered}, because inner blocks come first
in numbering sequence, not necessarily logical from a reader
perspective.

\end{fobx}

\end{fobx}

\phantomsection\label{custom-appearances}
\begin{fobx}{NewFeature}{\textbf{New feature 1.3: }Custom
\emph{appearances}}

yaml now allows for a new key group \texttt{appearances}. Here, you can
define general appearance of blocks and optional properties such as
color.

This way, appearance of multiple classes can be unified even if they do
not belong to the same numbering \texttt{group}

\phantomsection\label{change-in-color-specification}
\begin{fobx}{NewFeature}{\textbf{Change: }Change in color specification}

Colors should now be specified as hex strings starting with \texttt{\#}.
This is supported by most major editors, and the canonical way in quarto
\texttt{brand} (the old syntax is still working)

\end{fobx}

\end{fobx}

\phantomsection\label{bref}
\begin{fobx}{NewFeature}{\textbf{New feature 1.4: }\texttt{longref} for
cross references that include block name}

Cross references have previously only been possible in LaTeX style, via
\texttt{\textbackslash{}ref\{thelabel\}}. Now a new variant,
\texttt{\textbackslash{}longref\{thelabel\}}, has been added, that adds
the name of the environment as a reference prefix to the resolved link.
By default, the reference prefix is equal to the environment name, which
can be changed as attribute \texttt{label}. The reference prefix can be
redefined via the attribute \texttt{reflabel}.

\phantomsection\label{cnb-7-TODO}
\begin{fobx}{TODO}{\textbf{To do 1.1: }future plans for cross references}

\begin{itemize}
\tightlist
\item[$\square$]
  should I change \texttt{label} to \texttt{blockname}? It would be
  breaking for those who have used it.
\item[$\square$]
  consider making crossreferences using \texttt{@thelabel} as in plain
  quarto. Use the pandoc facilities for that
\end{itemize}

\end{fobx}

\end{fobx}

\phantomsection\label{custom-lists-of-blocks}
\begin{fobx}{Feature}{\textbf{Feature 1.5: }Custom lists-of blocks}

Generate \texttt{.qmd} files that contains a list of selected block
classes, intermitted with headers from the document for easy
customization and commenting.

This way, one can make a list of all definitions, examples, or
\{theorems, propositions and lemmas\} etc., edit it later and attach to
the main document. If you edit, make sure to rename the autogenerated
list first, otherwise it will get overwritten in the next run and all
work is lost \ldots{}

Currently, you need to give a key \texttt{listin} for any class or group
of classes that should appear in a list of things. The value of this key
is an array of names, also if only one list is to be generated. These
names are turned into files \texttt{list-of-}name\texttt{.qmd}. I am
considering replacing the yaml interface by a sub-key to
\texttt{custom-numbered-classes}. This would allow to define arbitrary
classes that can be attached to any custom div block, such as
\texttt{.important}.

\phantomsection\label{automatic-inclusion-of-lists-of}
\begin{fobx}{Tip}{\textbf{Tip: }automatic inclusion of lists-of}

You can include generated lists-of using a quarto shortcut. To do this,
you need to compile the document twice, in order to get the list-of
file.

\begin{tcolorbox}[enhanced jigsaw, coltitle=black, opacitybacktitle=0.6, colbacktitle=quarto-callout-warning-color!10!white, bottomrule=.15mm, left=2mm, arc=.35mm, colframe=quarto-callout-warning-color-frame, toptitle=1mm, colback=white, breakable, bottomtitle=1mm, opacityback=0, toprule=.15mm, title=\textcolor{quarto-callout-warning-color}{\faExclamationTriangle}\hspace{0.5em}{Warning}, titlerule=0mm, rightrule=.15mm, leftrule=.75mm]

You need to comment out the include shortcode until the first version of
the ``list-of'' qmd file is generated.

If your custom block has headers inside, it may break the div structure.
To avoid this, enclose the header n another div

\end{tcolorbox}

\paragraph*{New features mentioned in this
document:}\label{new-features-mentioned-in-this-document}
\addcontentsline{toc}{paragraph}{New features mentioned in this
document:}

\providecommand{\Pageref}[1]{\hfill p.\pageref{#1}}

\hyperref[cnb-1-NewFeature]{\textbf{New feature 1.1}}: New options for
block appearance\Pageref{cnb-1-NewFeature}

\hyperref[nested-blocks-may-be-numbered]{\textbf{New feature}}: Nested
blocks may be numbered\Pageref{nested-blocks-may-be-numbered}

\hyperref[change-in-color-specification]{\textbf{Change}}: Change in
color specification\Pageref{change-in-color-specification}

\hyperref[custom-appearances]{\textbf{New feature 1.3}}: Custom
\emph{appearances}\Pageref{custom-appearances}

\hyperref[bref]{\textbf{New feature 1.4}}: longref for cross references
that include block name\Pageref{bref}

\end{fobx}

\end{fobx}

\section{Pseudomath examples}\label{pseudomath-examples}

The blocks in this section are mostly \texttt{quartothmlike} blocks. The
section also demonstrates that the section prefix can be changed by
specifying the attribute \texttt{numberprefix} in the section heading.
So here, section number''2'' is replaced by ``P''.

\phantomsection\label{fancy}
\begin{fobx}{Definition}{\textbf{Definition P.1: }F\(\alpha\)ncybox}

A box is called \emph{f\(\alpha\)ncybox} if it looks quite fancy.

In this context, by \emph{fancy} we mean that the title of the box
appears as a clickable button when rendered as html, where
\emph{clickable} implies that it throws a small shadow that becomes
bigger when hovering over it.

\end{fobx}

\phantomsection\label{fancy-is-fancy}
\begin{fobx}{Corollary}{\textbf{Corollary P.2: }fancy is fancy}

By \hyperref[fancy]{Definition P.1}, \texttt{foldboxes} are fancyboxes.

\end{fobx}

\phantomsection\label{TeachersHope}
\begin{fobx}{Conjecture}{\textbf{Conjecture P.3}}

Students are lured into clicking on the title and unfolding the
fancybox.

\end{fobx}

\phantomsection\label{TeacherHopes}
\begin{fobx}{Theorem}{\textbf{Theorem P.4}}

This extension has been written by a teacher who hopes that students
read the course notes\ldots{}

\end{fobx}

\hyperref[TeacherHopes]{Theorem P.4} is suggested by Conjecture
\hyperref[TeachersHope]{P.3}, but it cannot be proven theoretically. It
does give rise to more conjectures, though. Here a reference \ref{Hope}
to an unknown object.

\phantomsection\label{cnb-14-Conjecture}
\begin{fobx}{Conjecture}{\textbf{Conjecture}}

The teacher mentioned in Theorem \hyperref[TeacherHopes]{P.4} is a
statistician who got addicted to quarto due to James J Balamuta's
fantastic \href{https://github.com/coatless/quarto-webr}{quarto-webr
extension}, and desparately wanted to have a common counter for theorems
and alike. She got also convinced that everything is possible in quarto
by the many nice extensions from Shafayet Khan Shafee.

\end{fobx}




\end{document}
